\documentclass[12pt]{article}
\usepackage[a4paper,margin=1in]{geometry}
\usepackage[utf8]{inputenc}
\usepackage[T1]{fontenc}
\usepackage{amsmath,amssymb}
\usepackage{textcomp}
\usepackage{parskip}
\begin{document}
\begin{center}
{\Large \textbf{2009 STEP II}}
\end{center}

\subsection*{Page 1}
\begin{flushleft}
STEP II, 2009\par
\bigskip\textbf{Section A: Pure Mathematics}\par
\bigskip\textbf{Question 1.} Two curves have equations x4 + y4 = u and xy = v , where u and v are positive constants.\par
State the equations of the lines of symmetry of each curve.\par
The curves intersect at the distinct points A, B, C and D (taken anticlockwise from A).\par
\[The coordinates of A are (α, β), where α > β > 0. Write down, in terms of α and β, the\]
coordinates of B, C and D.\par
Show that the quadrilateral ABCD is a rectangle and find its area in terms of u and v only.\par
\[Verify that, for the case u = 81 and v = 4, the area is 14.\]
\bigskip\textbf{Question 2.} The curve C has equation\par
\[y =\]
\[asin(\piex)\]
,\par
\[where a > 1.\]
\textbf{(i)} Find the coordinates of the stationary points on C.\par
t\par
\textbf{(ii)} Use the approximations e 1 + t and sin t t (both valid for small values of t) to show\par
\approx \approx\par
that\par
\[y 1 \pix ln a\]
\approx -\par
for small values of x.\par
\textbf{(iii)} Sketch C.\par
\textbf{(iv)} By approximating C by means of straight lines joining consecutive stationary points,\par
show that the area between C and the x-axis between the kth and (k + 1)th maxima is\par
approximately\par
\[a^2 + 1\]
ln 1 + k 3 ) -1 .\par
2a - 4\par
³ ´ ³ ´\par
¡\par
\end{flushleft}
\subsection*{Page 2}
\begin{flushleft}
\bigskip\textbf{Question 3.} Prove that\par
\[tan 1 \pi 1 x sec x tan x . ( )\]
4 - 2 \equiv - ∗\par
¡ ¢\par
\textbf{(i)} Use ( ) to find the value of tan 1 \pi . Hence show that\par
∗ 8\par
\[\sqrt\{3\} + \sqrt\{2\} 1\]
\[tan 11 \pi = - .\]
\[24 \sqrt\{3\} \sqrt\{6\} + 1\]
-\par
\textbf{(ii)} Show that\par
\[\sqrt\{3\} + \sqrt\{2\} 1\]
\[- = 2 + \sqrt\{2\} + \sqrt\{3\} + \sqrt\{6\} .\]
\[\sqrt\{3\} \sqrt\{6\} + 1\]
-\par
\textbf{(iii)} Use ( ) to show that\par
∗\par
\[tan 1 \pi = 16 + 10\sqrt\{2\} + 8\sqrt\{3\} + 6\sqrt\{6\} 2 \sqrt\{2\} \sqrt\{3\} \sqrt\{6\} .\]
48 - - - -\par
q\par
\bigskip\textbf{Question 4.} The polynomial p(x) is of degree 9 and p(x) 1 is exactly divisible by (x 1)5.\par
- -\par
\textbf{(i)} Find the value of p(1).\par
\textbf{(ii)} Show that p 0 (x) is exactly divisible by (x 1)4.\par
-\par
\textbf{(iii)} Given also that p(x) + 1 is exactly divisible by (x + 1)5, find p(x).\par
\[5 Expand and simplify (\sqrt\{x\} 1 + 1)2 .\]
-\par
\textbf{(i)} Evaluate\par
\[10 x + 2\sqrt\{x\} 1 + x 2\sqrt\{x\} 1\]
- - - dx .\par
\[\sqrt\{x\} 1\]
\[Z^5 p\]
-\par
p\par
\textbf{(ii)} Find the total area between the curve\par
\[x 2\sqrt\{x\} 1\]
\[y = - -\]
\[\sqrt\{x\} 1\]
p\par
-\par
\[and the x-axis between the points x = 5 and x = 10.\]
\textbf{(iii)} Evaluate\par
\[10 x + 2\sqrt\{x\} 1 + x 2\sqrt\{x\} + 1 + 2\]
- - dx .\par
\[Z 4 5 p \sqrt\{x^2\} p - 1\]
\end{flushleft}
\subsection*{Page 3}
\begin{flushleft}
\bigskip\textbf{Question 6.} The Fibonacci sequence F , F , F , . . . is defined by F = 1, F = 1 and\par
1 2 3 1 2\par
\[Fn+1 = Fn + Fn-1 (n > 2).\]
Write down the values of F , F , . . ., F .\par
3 4 10\par
\infty\par
\[Let S = .\]
Fi\par
\[i=1\]
X\par
1 1\par
\textbf{(i)} Show that > for i > 4 and deduce that S > 3 .\par
Fi 2Fi-1\par
\[Show also that S < 3 2 .\]
\textbf{(ii)} Show further that 3.2 < S < 3.5 .\par
\bigskip\textbf{Question 7.} Let y = (x a) n e bx\sqrt1 + x2 , where n and a are constants and b is a non-zero constant. Show\par
-\par
that\par
dy (x a)\par
n-1e bx\par
q(x)\par
\[= - ,\]
\[dx \sqrt\{1\} + x^2\]
where q(x) is a cubic polynomial.\par
Using this result, determine:\par
\[(x 4)14e4x (4x3 1)\]
\textbf{(i)} - - dx ;\par
\[\sqrt\{1\} + x^2\]
\[\int\]
\[(x 1)21e12x (12x4 x^2 11)\]
\textbf{(ii)} - - - dx ;\par
\[\sqrt\{1\} + x^2\]
\[\int\]
\[(x 2)6e4x (4x4 + x^3 2)\]
\textbf{(iii)} - - dx .\par
\[\sqrt\{1\} + x^2\]
\[\int\]
\bigskip\textbf{Question 8.} The non-collinear points A, B and C have position vectors a, b and c, respectively. The\par
points P and Q have position vectors p and q, respectively, given by\par
\[p = λa + (1 λ)b and q = µa + (1 µ)c\]
- -\par
\[where 0 < λ < 1 and µ > 1. Draw a diagram showing A, B, C, P and Q.\]
\[Given that CQ BP = AB AC, find µ in terms of λ, and show that, for all values of λ, the\]
\[\times \times\]
the line P Q passes through the fixed point D, with position vector d given by d = a + b + c .\par
-\par
What can be said about the quadrilateral ABDC?\par
\end{flushleft}
\subsection*{Page 4}
\begin{flushleft}
\bigskip\textbf{Section B: Mechanics}\par
\bigskip\textbf{Question 9.} (i) A uniform lamina OXY Z is in the shape of the trapezium shown in the diagram. It is\par
right-angled at O and Z, and OX is parallel to Y Z. The lengths of the sides are given\par
\[by OX = 9 cm, XY = 41 cm, Y Z = 18 cm and ZO = 40 cm. Show that its centre of\]
mass is a distance 7 cm from the edge OZ.\par
Z Y\par
40 41\par
O X\par
\textbf{(ii)} The diagram shows a tank with no lid made of thin sheet metal. The base OXU T , the\par
back OT W Z and the front XU V Y are rectangular, and each end is a trapezium as in\par
part (i) . The width of the tank is d cm.\par
W V\par
d\par
Z Y\par
40 41\par
T\par
U\par
d\par
O 9 X\par
Show that the centre of mass of the tank, when empty, is a distance\par
3(140 + 11d)\par
cm\par
5(12 + d)\par
from the back of the tank.\par
The tank is then filled with a liquid. The mass per unit volume of this liquid is k times\par
\[the mass per unit area of the sheet metal. In the case d = 20, find an expression for the\]
distance of the centre of mass of the filled tank from the back of the tank.\par
\end{flushleft}
\subsection*{Page 5}
\begin{flushleft}
y\par
u\par
O\par
x\par
P P\par
1 2\par
P\par
u\par
P\par
Four particles P , P , P and P , of masses m , m , m and m , respectively, are arranged on\par
1 2 3 4 1 2 3 4\par
smooth horizontal axes as shown in the diagram.\par
Initially, P and P are stationary, and both P and P are moving towards O with speed u.\par
2 3 1 4\par
Then P and P collide, at the same moment as P and P collide. Subsequently, P and P\par
1 2 4 3 2 3\par
collide at O, as do P and P some time later. The coefficient of restitution between each pair\par
1 4\par
\[of particles is e, and e > 0.\]
Show that initially P and P are equidistant from O.\par
2 3\par
\end{flushleft}
\subsection*{Page 6}
\begin{flushleft}
\bigskip\textbf{Question 11.} A train consists of an engine and n trucks. It is travelling along a straight horizontal section\par
of track. The mass of the engine and of each truck is M . The resistance to motion of the\par
engine and of each truck is R, which is constant. The maximum power at which the engine\par
can work is P .\par
Obtain an expression for the acceleration of the train when its speed is v and the engine is\par
working at maximum power.\par
The train starts from rest with the engine working at maximum power. Obtain an expression\par
for the time T taken to reach a given speed V , and show that this speed is only achievable if\par
\[P > (n + 1)RV .\]
\textbf{(i)} In the case when (n + 1)RV /P is small, use the approximation ln(1 x) x 1 x2\par
- \approx - - 2\par
(valid for small x) to obtain the approximation\par
P T 1 (n + 1)M V 2\par
\approx 2\par
and interpret this result.\par
\textbf{(ii)} In the general case, the distance moved from rest in time T is X. Write down, with\par
explanation, an equation relating P , T , X, M , V , R and n and hence show that\par
2P T (n + 1)M V 2\par
\[X = - .\]
2(n + 1)R\par
\end{flushleft}
\subsection*{Page 7}
\begin{flushleft}
\bigskip\textbf{Section C: Probability and Statistics}\par
\bigskip\textbf{Question 12.} A continuous random variable X has probability density function given by\par
\[0 for x < 0\]
\[f(x) =\]
(ke\par
\[-2x2\]
\[for 0 6 x < ,\]
\infty\par
where k is a constant.\par
\textbf{(i)} Sketch the graph of f(x).\par
\textbf{(ii)} Find the value of k.\par
\textbf{(iii)} Determine E(X) and Var (X).\par
\textbf{(iv)} Use statistical tables to find, to three significant figures, the median value of X.\par
13 Satellites are launched using two different types of rocket: the Andover and the Basingstoke.\par
The Andover has four engines and the Basingstoke has six. Each engine has a probability p\par
of failing during any given launch. After the launch, the rockets are retrieved and repaired by\par
replacing some or all of the engines. The cost of replacing each engine is K.\par
For the Andover, if more than one engine fails, all four engines are replaced. Otherwise, only\par
the failed engine (if there is one) is replaced. Show that the expected repair cost for a single\par
launch using the Andover is\par
\[4Kp(1 + q + q^2 2q3) (q = 1 p) (*)\]
- -\par
For the Basingstoke, if more than two engines fail, all six engines are replaced. Otherwise only\par
the failed engines (if there are any) are replaced. Find, in a form similar to ( ), the expected\par
∗\par
repair cost for a single launch using the Basingstoke.\par
Find the values of p for which the expected repair cost for the Andover is 2 of the expected\par
repair cost for the Basingstoke.\par
\end{flushleft}
\end{document}
